\subsection{Optimization Problem}
\label{subsec:problem_statement}

We formulate a broad class of smooth graph signal recovery problems as a unified constrained convex optimization problem incorporating the proposed regularizer to mitigate excessive penalties induced by an inaccurate proxy graph.

Let $\mathbf{u}^* \in \mathbb{R}^{|\mathcal{V}|}$ be the smooth graph signal to be recovered. We consider the following linear observation model: 
\begin{equation}
\mathbf{v}=\boldsymbol{\Phi}\mathbf{u}^*+\mathbf{n},
\label{eq:proposed_observation_model}
\end{equation}
where $\mathbf{v} \in \mathbb{R}^K$ is the observed signal, and $\mathbf{n} \in \mathbb{R}^K$ is additive Gaussian noise.
The linear observation matrix $\boldsymbol{\Phi}\in\mathbb{R}^{K\times|\mathcal{V}|}$ specifies the graph signal recovery task of interest, e.g., denoising when $\Phi = I$, inpainting when $\Phi$ is a binary masking matrix, and compressed sensing when $\Phi$ is a general sensing matrix.

Under this observation model, we formulate the recovery of the smooth graph signal $\mathbf{u}^*$ as the following constrained convex optimization problem using our proposed regularizer:
\begin{equation}
\begin{aligned}
\min_\mathbf{u}\quad
R_\mathrm{IC}(\mathbf{u})\quad
\mathrm{s.t.}\quad
&\left\{
\begin{aligned}
&\boldsymbol{\Phi}\mathbf{u}\in B_{\mathbf{v},r},\\
&\mathbf{u}\in \lbrack \underline{\mu},\overline{\mu} \rbrack^{|\mathcal{V}|},
\end{aligned}
\right.
\end{aligned}
\label{eq:proposed_constrained_problem}
\end{equation}
where $\mathbf{u}$ denotes the optimization variable corresponding to the recovered signal, and the closed convex constraint sets are defined as
\begin{equation}
B_{\mathbf{v},r}
:=\{\mathbf{z}\in\mathbb{R}^{K}\mid
\|\mathbf{z}-\mathbf{v}\|_2\leq r\},
\label{eq:l2_ball_constraint_set}
\end{equation}
with the hyperparameters $r > 0$, and $\underline{\mu}, \overline{\mu} \in \mathbb{R}$. 

The data-fidelity constraint $\boldsymbol{\Phi}\mathbf{u}\in B_{\mathbf{v},r}$ enforces the recovered signal and the observed signal to stay within a specified distance in the observation domain, i.e., $\lVert \boldsymbol{\Phi}\mathbf{u}-\mathbf{v}\rVert_2\leq r$.
Unlike a typical penalty formulation such as $\frac{\lambda}{2}\lVert \boldsymbol{\Phi}\mathbf{u}-\mathbf{v}\rVert_2^2$ with a balancing parameter $\lambda$, this constrained formulation directly specifies the admissible discrepancy between $\boldsymbol{\Phi}\mathbf{u}$ and $\mathbf{v}$ by $r$, not in a relative balance with other terms.
This independence allows a direct interpretation of data fidelity and facilitates principled hyperparameter control \cite{Afonso2010, OnoYamada2015, Chierchia2015, Ono2017PnP, Ono2017L0}.
In particular, $r$ can be determined from the standard deviation $\sigma$ of the Gaussian noise $\mathbf{n}$ as $r = \theta \sigma \sqrt{K}$, where $\theta>0$ is a practical adjustment parameter, typically chosen to be close to $1$.

The box constraint $\mathbf{u}\in \lbrack \underline{\mu},\overline{\mu} \rbrack^{|\mathcal{V}|}$ enforces the recovered signal to stay within the specified value range, i.e., $\underline{\mu} \leq u_i \leq \overline{\mu}$.
The upper and lower bounds, $\underline{\mu}$ and $\overline{\mu}$, can be set to the known range of signal values or sufficiently large when such prior knowledge is unavailable.
Additionally, the box constraint ensures Problem~\eqref{eq:proposed_constrained_problem} is well-posed in the sense that it admits at least one optimal solution, provided that the feasible set is nonempty.
Specifically, this constraint makes the feasible set $\mathcal{S}:=\lbrack \underline{\mu},\overline{\mu} \rbrack^{|\mathcal{V}|}\cap\boldsymbol{\Phi}^{-1}(B_{\mathbf{v},r})$ compact (closed and bounded), as $\mathcal{S}$ is a closed subset of the compact set $\lbrack \underline{\mu},\overline{\mu} \rbrack^{|\mathcal{V}|}$, which follows from the fact that $\boldsymbol{\Phi}^{-1}(B_{\mathbf{v},r}) = \{\mathbf{z} \mid \boldsymbol{\Phi} \mathbf{z} \in B_{\mathbf{v}, r}\}$ is closed.
Thus, the lower semicontinuous objective function $R_\mathrm{IC}$ attains its infimum over the feasible set.


\subsection{P-PDS-Based Optimization Algorithm}
\label{subsec:optimization_algorithm}

We develop an efficient algorithm for solving
Problem~\eqref{eq:proposed_constrained_problem} based on P-PDS, which
requires neither matrix inversions nor inner iterative solvers and allows
convergence-guaranteed variable-wise stepsize selection.
Using
$R_{\mathrm{IC}}(\mathbf{u})
=
\Psi_{\tau}(\sqrt{\mathbf{W}}\mathbf{D}\mathbf{u})$
and the indicator functions of the constraint sets,
Problem~\eqref{eq:proposed_constrained_problem} can be equivalently
reformulated as
\begin{equation}
\begin{aligned}
\min_{\mathbf{u},\mathbf{y}_1,\mathbf{y}_2}\quad
&
\iota_{B_{\mathbf{v},r}}(\mathbf{y}_1)
+
\iota_{\lbrack\underline{\mu},\overline{\mu}\rbrack^{|\mathcal{V}|}}
(\mathbf{u})
+
\Psi_{\tau}(\mathbf{y}_2)
\\
\mathrm{s.t.}\quad
&
\left\{
\begin{aligned}
&\mathbf{y}_1=\boldsymbol{\Phi}\mathbf{u},\\
&\mathbf{y}_2=\sqrt{\mathbf{W}}\mathbf{D}\mathbf{u}.
\end{aligned}
\right.
\end{aligned}
\label{eq:p_pds_formulation}
\end{equation}
This formulation fits the canonical P-PDS problem in
\eqref{eq:p_pds_general_form} by setting
$\mathbf{x}_1=\mathbf{u}$,
$f_1=
\iota_{\lbrack\underline{\mu},\overline{\mu}\rbrack^{|\mathcal{V}|}}$,
$g_1=\iota_{B_{\mathbf{v},r}}$, and
$g_2=\Psi_{\tau}$, with
$\mathbf{L}_{1,1}=\boldsymbol{\Phi}$ and
$\mathbf{L}_{2,1}=\sqrt{\mathbf{W}}\mathbf{D}$.

The required proximity operators are available in closed form.
For $\kappa>0$, the proximity operators of the two indicator functions
are given by
\begin{equation}
\begin{aligned}
&\operatorname{prox}_{
\kappa
\iota_{\lbrack\underline{\mu},\overline{\mu}\rbrack^{|\mathcal{V}|}}
}
(\mathbf{x})
=
\max\!\left(
\underline{\mu},
\min\!\left(\overline{\mu},\mathbf{x}\right)
\right),
\\
&\operatorname{prox}_{\kappa\iota_{B_{\mathbf{v},r}}}(\mathbf{x})
=
\begin{cases}
\mathbf{x},
& \|\mathbf{x}-\mathbf{v}\|_2\leq r,
\\[1mm]
\displaystyle
\mathbf{v}
+
\frac{r(\mathbf{x}-\mathbf{v})}
{\|\mathbf{x}-\mathbf{v}\|_2},
& \text{otherwise},
\end{cases}
\end{aligned}
\label{eq:prox_constraint_sets}
\end{equation}
where the maximum and minimum are taken componentwise.
Moreover, as noted in Remark~\ref{rem:proximability}, $\Psi_{\tau}$ is proximable, and its proximity operator is given by
\begin{equation}
\operatorname{prox}_{\kappa\Psi_{\tau}}(\mathbf{x})
=
\frac{1}{1+2\kappa}\left(\mathbf{x}
+
2\kappa\operatorname{soft}
\left(
\mathbf{x};
(1+2\kappa)\tau
\right)\right).
\label{eq:prox_psi}
\end{equation}
Hence, all proximity operations required by P-PDS can be evaluated
explicitly.
The resulting algorithm is summarized in
Algorithm~\ref{alg:proposed_algorithm}.

\begin{algorithm}[t]
\caption{P-PDS-based algorithm for solving
Problem~\eqref{eq:proposed_constrained_problem}}
\label{alg:proposed_algorithm}
\begin{algorithmic}[1]
\STATE \textbf{Input:}
$\mathbf{u}^{(0)}$,
$\mathbf{z}_1^{(0)}$,
$\mathbf{z}_2^{(0)}$
\STATE $n\leftarrow 0$
\WHILE{Stopping criteria are not met}
\STATE
$\widetilde{\mathbf{u}}
\leftarrow
\boldsymbol{\Phi}^{\top}\mathbf{z}_1^{(n)}
+
(\sqrt{\mathbf{W}}\mathbf{D})^{\top}\mathbf{z}_2^{(n)}$
\STATE
$\mathbf{u}^{(n+1)}
\leftarrow
\operatorname{prox}_{
\rho_1
\iota_{\lbrack\underline{\mu},\overline{\mu}\rbrack^{|\mathcal{V}|}}
}
\left(
\mathbf{u}^{(n)}
-
\rho_1\widetilde{\mathbf{u}}
\right)$
\STATE
$\widetilde{\mathbf{z}}_1
\leftarrow
\mathbf{z}_1^{(n)}
+
\gamma_1
\boldsymbol{\Phi}
\left(
2\mathbf{u}^{(n+1)}
-
\mathbf{u}^{(n)}
\right)$
\STATE
$\mathbf{z}_1^{(n+1)}
\leftarrow
\widetilde{\mathbf{z}}_1
-
\gamma_1
\operatorname{prox}_{\gamma_1^{-1}\iota_{B_{\mathbf{v},r}}}
\left(
\gamma_1^{-1}\widetilde{\mathbf{z}}_1
\right)$
\STATE
$\widetilde{\mathbf{z}}_2
\leftarrow
\mathbf{z}_2^{(n)}
+
\gamma_2
\sqrt{\mathbf{W}}\mathbf{D}
\left(
2\mathbf{u}^{(n+1)}
-
\mathbf{u}^{(n)}
\right)$
\STATE
$\mathbf{z}_2^{(n+1)}
\leftarrow
\widetilde{\mathbf{z}}_2
-
\gamma_2
\operatorname{prox}_{\gamma_2^{-1}\Psi_{\tau}}
\left(
\gamma_2^{-1}\widetilde{\mathbf{z}}_2
\right)$
\STATE $n\leftarrow n+1$
\ENDWHILE
\STATE \textbf{Output:} $\mathbf{u}^{(n)}$
\end{algorithmic}
\end{algorithm}

To determine the stepsizes, we employ the second variable-wise stepsize design proposed in \cite{NaganumaOno2023}, referred to therein as OVDP2.
For the matrices in \eqref{eq:p_pds_formulation}, let $\nu_{j,1}$ denote
upper bounds satisfying
$\|\mathbf{L}_{j,1}\|_{\mathrm{op}}\leq\nu_{j,1}$.
For $\mathbf{L}_{1,1}=\boldsymbol{\Phi}$, any upper bound
$\nu_{1,1}\geq\|\boldsymbol{\Phi}\|_{\mathrm{op}}$ can be used, and the
exact operator norm can be obtained as the largest singular value of
$\boldsymbol{\Phi}$.
In particular, $\nu_{1,1}=1$ can be used when
$\boldsymbol{\Phi}=\mathbf{I}$ or when $\boldsymbol{\Phi}$ is a binary
masking matrix.
For $\mathbf{L}_{2,1}=\sqrt{\mathbf{W}}\mathbf{D}$, we take the exact
operator norm
$\nu_{2,1}
=\|\sqrt{\mathbf{W}}\mathbf{D}\|_{\mathrm{op}}
=\sqrt{\lambda_{\max}(\mathbf{L})}$,
as shown in the proof of
Proposition~\ref{prop:mathematical-properties}.
Applying OVDP2 to these bounds yields
\begin{equation}
\begin{aligned}
\rho_1
&=
\frac{1}
{\nu_{1,1}+\sqrt{\lambda_{\max}(\mathbf{L})}},
\\
\gamma_1
&=
\frac{1}{\nu_{1,1}},
\qquad
\gamma_2
=
\frac{1}{\sqrt{\lambda_{\max}(\mathbf{L})}}.
\end{aligned}
\label{eq:stepsizes_proposed}
\end{equation}
Under the standard qualification condition stated in Section~II-C, the resulting sequence converges to a solution of Problem~\eqref{eq:proposed_constrained_problem}.